% GENERATED FILE. Edit canonical lessons, then run npm run generate:books.
% canonical-source-hash: fnv1a64:9ee92273d7eab7d3
% canonical-lessons: GE-C265-inzwischen, GE-C265-mittlerweile, GE-C265-waehrenddessen, GE-C265-bisher, GE-C265-laengst

\chapter{Inzwischen}
\label{ch:ge-inzwischen}
\hlchaptermodality{265}

\begin{chapteropening}
\textbf{By the end of this chapter:} \emph{I can say what happened in the meantime and up to now: meanwhile, by now, in the meantime, so far and long ago.}

Complete the last lesson of chapter 265: I can say what happened in the meantime and up to now: meanwhile, by now, in the meantime, so far and long ago.
\end{chapteropening}

\section[inzwischen]{inzwischen — meanwhile; by now}

\label{lesson:GE-C265-inzwischen}



\begin{quote}
Before the new one: say the German for recently, then the German for the day before yesterday.
\end{quote}

\subsection*{You'll want to know: inzwischen}
\textbf{inzwischen} — \textquotedblleft{}meanwhile; by now\textquotedblright{}.

\begin{quote}
Ich habe gekocht, und inzwischen hat er die Teller geholt.
\end{quote}

\begin{quote}
\emph{I cooked, and meanwhile he fetched the plates.}
\end{quote}

Also 'by now': \emph{Inzwischen ist sie zehn Jahre alt}.

\begin{cousinweb}
\emph{inzwischen} joins \emph{in} and \emph{zwischen} \textquotedblleft{}between\textquotedblright{}. \emph{Zwischen} comes from Old High German \emph{in zuiskēn} \textquotedblleft{}between the two\textquotedblright{}, built on \emph{zuiski} \textquotedblleft{}twofold\textquotedblright{} and so on the word for \textquotedblleft{}two\textquotedblright{} (EtymDB).
\end{cousinweb}

\begin{grammarlens}[title={What you've built}]
One more word for telling, explaining and describing at B1.
\end{grammarlens}

\subsection*{Your turn}
\begin{itemize}
\raggedright
  \item \emph{Say it:} \emph{inzwischen}
  \item \emph{Say it:} \emph{inzwischen}, once more
  \item \emph{Say it:} a sentence of your own with \emph{inzwischen}
  \item \emph{Recall:} say the German for recently, then the German for the day before yesterday, then say \emph{inzwischen} again
\end{itemize}

\begin{tcolorbox}[breakable,colback=teal!4,colframe=teal!35!black,title={Before you move on}]
What does inzwischen mean? (\textquotedblleft{}Meanwhile; by now\textquotedblright{}.) Say it once more.
\end{tcolorbox}
\section[mittlerweile]{mittlerweile — by now, in the meanwhile}

\label{lesson:GE-C265-mittlerweile}



\begin{quote}
Before the new one: say the German for the day before yesterday, then the German for meanwhile.
\end{quote}

\subsection*{You'll want to know: mittlerweile}
\textbf{mittlerweile} — \textquotedblleft{}by now, in the meanwhile\textquotedblright{}.

\begin{quote}
Mittlerweile wohnt sie in einer größeren Wohnung.
\end{quote}

\begin{quote}
\emph{By now she lives in a bigger flat.}
\end{quote}

Close to \emph{inzwischen}, and a little more conversational.

\begin{cousinweb}
\emph{mittlerweile} is built on \emph{mittler} \textquotedblleft{}middle\textquotedblright{} and \emph{Weile} \textquotedblleft{}a while\textquotedblright{}, both inherited Germanic words; \emph{Weile} is a cousin of English \emph{while}. EtymDB lists no parent for the compound.
\end{cousinweb}

\begin{grammarlens}[title={What you've built}]
One more word for telling, explaining and describing at B1.
\end{grammarlens}

\subsection*{Your turn}
\begin{itemize}
\raggedright
  \item \emph{Say it:} \emph{mittlerweile}
  \item \emph{Say it:} \emph{mittlerweile}, once more
  \item \emph{Say it:} a sentence of your own with \emph{mittlerweile}
  \item \emph{Recall:} say the German for the day before yesterday, then the German for meanwhile, then say \emph{mittlerweile} again
\end{itemize}

\begin{tcolorbox}[breakable,colback=teal!4,colframe=teal!35!black,title={Before you move on}]
What does mittlerweile mean? (\textquotedblleft{}By now, in the meanwhile\textquotedblright{}.) Say it once more.
\end{tcolorbox}
\section[währenddessen]{währenddessen — in the meantime, while that was going on}

\label{lesson:GE-C265-waehrenddessen}



\begin{quote}
Before the new one: say the German for meanwhile, then the German for by now.
\end{quote}

\subsection*{You'll want to know: währenddessen}
\textbf{währenddessen} — \textquotedblleft{}in the meantime, while that was going on\textquotedblright{}.

\begin{quote}
Sie telefonierte. Währenddessen kochte das Wasser.
\end{quote}

\begin{quote}
\emph{She was on the phone. In the meantime the water boiled.}
\end{quote}

Built on \emph{während}. As an adverb, the verb comes second after it.

\begin{cousinweb}
\emph{währenddessen} joins \emph{während} \textquotedblleft{}during\textquotedblright{}, the present participle of \emph{währen} \textquotedblleft{}to last\textquotedblright{} (EtymDB), and \emph{dessen} \textquotedblleft{}of that\textquotedblright{}: while that lasted.
\end{cousinweb}

\begin{grammarlens}[title={What you've built}]
One more word for telling, explaining and describing at B1.
\end{grammarlens}

\subsection*{Your turn}
\begin{itemize}
\raggedright
  \item \emph{Say it:} \emph{währenddessen}
  \item \emph{Say it:} \emph{währenddessen}, once more
  \item \emph{Say it:} a sentence of your own with \emph{währenddessen}
  \item \emph{Recall:} say the German for meanwhile, then the German for by now, then say \emph{währenddessen} again
\end{itemize}

\begin{tcolorbox}[breakable,colback=teal!4,colframe=teal!35!black,title={Before you move on}]
What does währenddessen mean? (\textquotedblleft{}In the meantime, while that was going on\textquotedblright{}.) Say it once more.
\end{tcolorbox}
\section[bisher]{bisher — so far, until now}

\label{lesson:GE-C265-bisher}



\begin{quote}
Before the new one: say the German for by now, then the German for in the meantime.
\end{quote}

\subsection*{You'll want to know: bisher}
\textbf{bisher} — \textquotedblleft{}so far, until now\textquotedblright{}.

\begin{quote}
Bisher ist alles gut gegangen.
\end{quote}

\begin{quote}
\emph{So far everything has gone well.}
\end{quote}

With the present or the \emph{Perfekt}: \emph{Bisher habe ich nichts gehört}.

\begin{cousinweb}
\emph{bisher} joins \emph{bis} \textquotedblleft{}until\textquotedblright{} and \emph{her} \textquotedblleft{}hither, to here\textquotedblright{}, so \textquotedblleft{}up to here\textquotedblright{}. Both are inherited Germanic words; EtymDB lists no parent.
\end{cousinweb}

\begin{grammarlens}[title={What you've built}]
One more word for telling, explaining and describing at B1.
\end{grammarlens}

\subsection*{Your turn}
\begin{itemize}
\raggedright
  \item \emph{Say it:} \emph{bisher}
  \item \emph{Say it:} \emph{bisher}, once more
  \item \emph{Say it:} a sentence of your own with \emph{bisher}
  \item \emph{Recall:} say the German for by now, then the German for in the meantime, then say \emph{bisher} again
\end{itemize}

\begin{tcolorbox}[breakable,colback=teal!4,colframe=teal!35!black,title={Before you move on}]
What does bisher mean? (\textquotedblleft{}So far, until now\textquotedblright{}.) Say it once more.
\end{tcolorbox}
\section[längst]{längst — long ago, long since}

\label{lesson:GE-C265-laengst}



\begin{quote}
Before the new one: say the German for in the meantime, then the German for so far.
\end{quote}

\subsection*{You'll want to know: längst}
\textbf{längst} — \textquotedblleft{}long ago, long since\textquotedblright{}.

\begin{quote}
Der Zug ist längst abgefahren.
\end{quote}

\begin{quote}
\emph{The train left long ago.}
\end{quote}

\emph{Längst nicht} means 'not nearly': \emph{Das ist längst nicht alles}.

\begin{cousinweb}
\emph{längst} is the superlative of \emph{lang} \textquotedblleft{}long\textquotedblright{}, an inherited Germanic word (EtymDB): longest ago.
\end{cousinweb}

\begin{grammarlens}[title={What you've built}]
One more word for telling, explaining and describing at B1.
\end{grammarlens}

\subsection*{Your turn}
\begin{itemize}
\raggedright
  \item \emph{Say it:} \emph{längst}
  \item \emph{Say it:} \emph{längst}, once more
  \item \emph{Say it:} a sentence of your own with \emph{längst}
  \item \emph{Recall:} say the German for in the meantime, then the German for so far, then say \emph{längst} again
\end{itemize}

\begin{tcolorbox}[breakable,colback=teal!4,colframe=teal!35!black,title={Before you move on}]
What does längst mean? (\textquotedblleft{}Long ago, long since\textquotedblright{}.) Say it once more.
\end{tcolorbox}
